\documentclass{article}
\usepackage[a4paper, total={6.5in, 10in}]{geometry}
\setlength{\parindent}{0pt}
\usepackage{tcolorbox}
\tcbset{colback=white!94!black, colframe=black!60!white, coltitle = white, boxrule=0.8pt, arc=2mm}
\newtcolorbox{boxs}[2][]{title= #2,#1}
\usepackage{datetime2}
\usepackage[british]{babel}
\usepackage{amsfonts}
\usepackage{amsmath}

\title{\textbf{Kinematical Equations}}
\author{Robert O'Connell}
\date{29th September 2025}
\begin{document}
\maketitle

\section*{Kinematic Vector Definitions}

Each of the kinematic quantities (position, displacement, velocity, etc.) are defined as follows:

\subsection*{Position Vector}

The position vector at time $t$ is given by:
\[
\vec{r}(t) = \begin{bmatrix} x(t) \\ y(t) \end{bmatrix}
\]

\subsection*{Displacement Vector}

The displacement vector from time $t_1$ to time $t_2$ is:
\begin{align*}
\vec{s}(t_1, t_2) &= \vec{r}(t_2) - \vec{r}(t_1) \\
&= \begin{bmatrix} x(t_2) - x(t_1) \\ y(t_2) - y(t_1) \end{bmatrix}
\end{align*}

\subsection*{Velocity Vector}

The instantaneous velocity at time $t$ is:
\begin{align*}
\vec{v}(t) &= \lim_{\Delta t\to0}\frac{\vec{r}(t + \Delta t) - \vec{r}(t)}{\Delta t} \\
&= \frac{d\vec{r}}{dt}
\end{align*}

\subsection*{Acceleration Vector}

The instantaneous acceleration at time $t$ is:
\begin{align*}
\vec{a}(t) &= \lim_{\Delta t\to0}\frac{\vec{v}(t + \Delta t) - \vec{v}(t)}{\Delta t} \\
&= \frac{d\vec{v}}{dt}
\end{align*}

\subsection*{Component-wise Representation}

The position vector can be written in component form as:
\[
\vec{r}(t) = x(t)\hat{i} + y(t)\hat{j}
\]

By differentiating with respect to time:
\begin{align*}
    \vec{v}(t) &= \frac{dx(t)}{dt}\hat{i} + x(t)\frac{d\hat{i}}{dt} + \frac{dy(t)}{dt}\hat{j} + y(t)\frac{d\hat{j}}{dt} \\
    &= \frac{dx(t)}{dt}\hat{i} + \frac{dy(t)}{dt}\hat{j}
\end{align*}

\noindent
Note that $\frac{d\hat{i}}{dt} = \frac{d\hat{j}}{dt} = 0$ since the Cartesian basis vectors are constant (independent of time).

\noindent
The acceleration can similarly be decomposed:
\[
\vec{a}(t) = \frac{d^2x(t)}{dt^2}\hat{i} + \frac{d^2y(t)}{dt^2}\hat{j}
\]

\noindent
These equations can also be extended to three dimensions by adding a $z$-component with basis vector $\hat{k}$.

\newpage

\section*{Formal Solution of Kinematic Equations}

Suppose $\vec{a}(t)$ is known. How do we determine $\vec{r}(t)$ and $\vec{v}(t)$?

\noindent
In three dimensions, the acceleration vector is:
\begin{align*}
    \vec{a}(t) &= \frac{dv_x(t)}{dt}\hat{i} + \frac{dv_y(t)}{dt}\hat{j} + \frac{dv_z(t)}{dt}\hat{k} \\ 
    &= a_x(t)\hat{i} + a_y(t)\hat{j} + a_z(t)\hat{k}
\end{align*}

\noindent
Focusing on the $x$-component:
\begin{align*}
    \frac{dv_x(t)}{dt} &= a_x(t) \\
    \int_{v_x(t_0)}^{v_x(t)} dv_x &= \int_{t_0}^{t} a_x(t')\,dt' \\
    v_x(t) - v_x(t_0) &= \int_{t_0}^{t} a_x(t')\,dt' \\
    v_x(t) &= v_x(t_0) + \int_{t_0}^{t} a_x(t')\,dt'
\end{align*}

\noindent
Returning to three dimensions, the velocity vector is:
\[
\vec{v}(t) = \vec{v}(t_0) + \int_{t_0}^{t} \vec{a}(t')\,dt'
\]

\noindent
Similarly, we can determine the position vector. Starting from $\frac{d\vec{r}}{dt} = \vec{v}(t)$:
\[
\vec{r}(t) = \vec{r}(t_0) + \int_{t_0}^{t} \vec{v}(t')\,dt'
\]

\noindent
These integral equations allow us to find the velocity and position given the acceleration and appropriate initial conditions $\vec{v}(t_0)$ and $\vec{r}(t_0)$.

\end{document}